% =====================================================================
%  Appendix -- Reproducibility
% =====================================================================

\section{Reproducibility}
\label{sec:appendix-repro}

\footnotesize
\sloppy
\setlength{\emergencystretch}{3em}

We provide the following artifacts for reproducibility.

\paragraph{Artifacts.}
\begin{itemize}\setlength\itemsep{2pt}
\item \textbf{Code repository} (MIT licence) at \repoURL; the \texttt{v1.0-ijdrr} tag is the camera-ready commit. Zenodo DOI: \texttt{[to be minted at submission]}.
\item \textbf{Pre-trained checkpoints} (\texttt{*.pt} PyTorch state-dicts) under \texttt{code/runs/3fold/}: 14 configurations $\times$ 3 folds $= 42$ checkpoints, $\sim 700$\,MB total.
\item \textbf{ONNX exports} (FP32 and INT8, per fold) for teacher and three students under \texttt{code/exports/}; $\sim 150$\,MB total.
\item \textbf{Per-fold metrics} in \texttt{results.json} (FP32) and \texttt{quantized.json} (INT8) under \texttt{code/runs/3fold/}; re-derivable by \texttt{run\_full\_experiment.py}.
\item \textbf{Figures} \ref{fig:f6-architecture}, \ref{fig:f7-detections} and \ref{fig:f8-inference} and the graphical abstract are generated by \texttt{make\_architecture\_figure.py}, \texttt{make\_detection\_gallery.py}, \texttt{make\_inference\_flow.py} and \texttt{make\_graphical\_abstract.py} from the released ONNX exports and result files.
\item \textbf{Per-platform latency} in \texttt{lat\_m2.json} and \texttt{lat\_wasm.json}; re-derivable by \texttt{bench\_m2.py} and \texttt{bench.mjs} respectively. Second-host latency (Table~\ref{tab:t13-host2}) in \texttt{code/runs/host2/\{lat\_ort\_cpu,lat\_wasm\_node,lat\_ort\_cpu\_u8s8\}.json}; re-derivable by \texttt{bench\_ort\_cpu.py} and \texttt{bench\_wasm\_node.mjs}.
\item \textbf{Quantization ablation} (Table~\ref{tab:t12-ptq-mechanism}) in \texttt{ptq\_mechanism.json} and the MobileViT layer-group diagnostic in \texttt{ptq\_mobilevit\_groups.json} under \texttt{code/runs/3fold/}; re-derivable by \texttt{ptq\_mechanism\_ablation.py} and \texttt{ptq\_mobilevit\_diagnostic.py}. The fold-0 unsigned-activation check is \texttt{code/runs/host2/u8s8\_fold0.json}.
\item \textbf{Zero-shot evaluations} in \texttt{ood\_fold\{0,1,2\}.json} (Sen1Floods11) and \texttt{zero\_shot\_fold\{0,1,2\}.json} (AIFloodSense $+$ FloodNet test, per-image metrics incl.\ flood-only IoU and false-alarm rates); re-derivable by \texttt{eval\_ood.py} and \texttt{eval\_zero\_shot.py}. The per-channel INT8 rows of Table~\ref{tab:t9-zeroshot} are the \texttt{*\_onnx\_int8\_pc} entries of the same files, produced by \texttt{eval\_zero\_shot\_onnx.py} from the per-channel artifacts of \texttt{ptq\_mechanism\_ablation.py}; \texttt{aggregate\_zero\_shot.py} builds \texttt{zero\_shot\_3fold.json}. The FSSD$\leftrightarrow$combined duplicate-check report is \texttt{fssd\_overlap\_report.json}, re-derivable by \texttt{check\_fssd\_overlap.py}.
\item \textbf{Combined-corpus retraining} (Section~\ref{sec:results-combined}) under \texttt{code/runs/combined/}: per-fold \texttt{results\_fold\{0,1,2\}.json} (re-derivable by \texttt{run\_full\_experiment.py --dataset combined} or \texttt{notebooks/floodlite\_combined\_kaggle.ipynb}), the three-fold aggregate \texttt{summary.json} (\texttt{aggregate\_combined.py}), the per-image KD contrasts and teacher-inheritance analysis \texttt{stats\_combined.json} (\texttt{stats\_combined.py}), and the INT8 results \texttt{quantized.json} (per-tensor; \texttt{quantize\_combined.py}) and \texttt{quantized\_pc.json} (per-channel; \texttt{quantize\_combined\_pc.py}). Checkpoints: 7 configurations $\times$ 3 folds $= 21$, $\sim 360$\,MB.
\end{itemize}

\paragraph{Environment.}
\begin{itemize}\setlength\itemsep{2pt}
\item Random seed: 42, fixed in the \texttt{KFold} constructor (data splits) and the notebook preamble; training-time stochasticity (weight init, augmentation order) is not pinned and is reflected in cross-fold standard deviations.
\item Software: Python 3.11; PyTorch 2.8 with CUDA 12.1 (Kaggle T4) or Apple-Silicon MPS (M2); \texttt{segmentation-models-pytorch} $\geq 0.4$; \texttt{timm} 1.0; \texttt{onnxruntime} 1.22.1 (both CPUs; INT8 QDQ artifacts require $\geq 1.18$); \texttt{onnxruntime-web} 1.17.3 (browser, M2 host) and 1.22.0 (browser, x86 host). The quantization ablation ran under PyTorch 2.15 (CPU), \texttt{onnx} 1.17, \texttt{timm} 1.0.15.
\item Hardware: training on a single Kaggle Tesla T4 (16\,GB); quantization and CPU latency on an Apple M2~Pro (10-core CPU, 16\,GB unified memory, macOS 26.2), WASM latency via Node.js~25; quantization ablation and second-host latency on an AMD Ryzen~5~5600G (6 cores, 16\,GB, Windows~10), WASM via Node.js~22.
\item Estimated reproduction cost: $\sim 35$ Kaggle-T4 hours $+ \sim 2$ M2-CPU hours for the FSSD benchmark, plus $\sim 30$ Kaggle-T4 hours (three sessions of under 12\,h) for the combined-corpus retraining; wall-clock $\approx 3$ calendar days assuming no Kaggle queue.
\end{itemize}
